\chapter{Conclusion}

\section{Result}


\section{Application Of The Result}
\subsubsection*{LMC and SMC Distances}
To evaluate the effectiveness of the calibration, distances to the Large and Small Magellanic Clouds (LMC and SMC) are derived using 24 versions of the calibrated Leavitt Law, covering four photometric bands (VIJK) and six color-based Wesenheit functions. This approach allows the construction of a distance matrix that highlights the relative performance of different color indices in producing consistent distance estimates. The results demonstrate the sensitivity of the PL calibration to the choice of photometric band and Wesenheit function, providing guidance for selecting the most robust combination for extragalactic distance determinations.



\subsection{Understanding the Universe: 1990s}

A major milestone in the extragalactic distance ladder was achieved in 1998 with the incorporation of Type~Ia supernovae as standardizable candles. The development of robust empirical methods to measure distances to galaxies hosting these supernovae, calibrated using the Cepheid period--luminosity relation, extended reliable distance measurements from scales of $\sim 100$~million light-years to several billion light-years. This dramatic extension of the cosmic distance scale led to one of the most surprising discoveries in modern cosmology: the expansion of the Universe is accelerating.

This acceleration implies the presence of a component with negative effective pressure that drives the expansion of spacetime. This unknown component was termed \emph{dark energy}. Despite extensive observational and theoretical efforts, the physical nature of dark energy remains unknown. Its discovery has underscored the importance of understanding not only the geometry of spacetime but also the detailed composition of the Universe. A simplified description of the Universe can be framed in terms of its total mass--energy content and the geometry it induces, where energy appears in various forms such as mass, radiation, and thermal energy.

\input{chapters/content/0_1_plot_uni_composition.tex}

In the standard cosmological framework, the Universe is composed primarily of three components: baryonic matter, dark matter, and dark energy. Ordinary, or \emph{baryonic}, matter consists of protons, neutrons, and electrons and accounts for only about $5\%$ of the total energy density of the Universe. This component forms the stars, galaxies, and large-scale structures directly observable through electromagnetic radiation.

Dark matter constitutes approximately $27\%$ of the total energy content of the Universe. It is hypothesized to possess fundamentally different properties from baryonic matter, as it neither emits nor absorbs electromagnetic radiation, rendering it invisible to direct observation. Its existence is inferred through its gravitational influence on visible matter and radiation. Strong observational evidence for dark matter arises from gravitational lensing, galaxy rotation curves, and the fine-scale anisotropies observed in the cosmic microwave background (CMB). Importantly, the interpretation of these phenomena depends critically on accurate distance measurements; for example, lensing geometries and dynamical mass estimates are highly sensitive to the assumed distances of the relevant systems. Despite its well-established gravitational effects, the physical nature of dark matter remains one of the most significant open questions in contemporary astrophysics.

Dark energy dominates the cosmic energy budget, contributing approximately $68\%$ of the total energy density of the Universe. In contrast, all other components---including radiation, neutrinos, and possible exotic particles---together contribute less than $1\%$. Determining the relative proportions of these components relies heavily on precise distance measurements across cosmic scales, as such measurements underpin estimates of luminosity, redshift, and volume, all of which are essential for reconstructing the Universe’s composition and its evolutionary history.

\subsubsection{$\Lambda$CDM Cosmological Model}

Observations of the distant Universe allow us to probe earlier cosmic epochs, as light requires finite time to traverse space. This fundamental property enables direct investigation of the Universe’s evolution. One approach to characterizing cosmic expansion involves measuring distances to the most distant observable objects, effectively constraining the size of the observable Universe and the cosmological horizon.

By combining insights from high-energy particle physics with general relativity, cosmologists have been able to model the evolution of the Universe back to fractions of a second after the Big Bang. Integrating these theoretical frameworks with a broad range of observational data has led to the formulation of the standard cosmological model, known as the $\Lambda$CDM (Lambda Cold Dark Matter) model. This model is characterized by a small set of parameters, including the densities of dark energy, dark matter, and baryonic matter; the Hubble constant ($H_0$); the scalar spectral index describing primordial density fluctuations; and the curvature parameter that determines the geometry of spacetime. Many of these parameters are tightly constrained by measurements of the cosmic microwave background and are independently supported by large-scale structure surveys, baryon acoustic oscillations, and supernova observations.

The full-sky temperature anisotropy map shown on the introductory page is a snapshot of the early Universe obtained by the \emph{Planck} satellite, a mission of the European Space Agency operational between 2009 and 2013. This map depicts the cosmic microwave background radiation, the relic thermal emission from the epoch of recombination. The small temperature fluctuations observed across the sky correspond to primordial density perturbations in the early Universe. Over billions of years, these perturbations grew through gravitational instability, ultimately giving rise to galaxies, clusters, and the large-scale structure observed today.

The CMB represents the earliest observable relic of the Universe, providing a view of conditions approximately 380{,}000~years after the Big Bang. Its precise measurement enables robust constraints on fundamental cosmological parameters, including the age, composition, and expansion rate of the Universe. Among the most significant results derived from the \emph{Planck} mission is a high-precision determination of the Hubble constant, $H_0$, which plays a central role in understanding both the past evolution and the future fate of the Universe.

\subsection{Cosmic Distance Ladder and the Hubble Constant: The Twenty-First Century}

Distance measurements across cosmic scales rely on a hierarchical set of interconnected methods collectively known as the \emph{cosmic distance ladder}, in which each rung is calibrated using more local and direct techniques. This framework allows distances to be extended progressively to cosmological scales, where they can be compared with early-Universe inferences from the cosmic microwave background (CMB).

\input{chapters/content/0_1_plot_ladder.tex}

A major refinement of the ladder was achieved by Riess and collaborators using Type~Ia supernovae (SNIa) as distance indicators. These thermonuclear explosions of white dwarfs exhibit a well-defined correlation between peak luminosity and light-curve shape, enabling their use as standardizable candles. By calibrating SNIa luminosities with Cepheid variable stars via the Leavitt Law, distances were extended from the Cepheid-accessible local volume of $\sim 30$~Mpc ($\sim 10^8$~ly) to beyond $1000$~Mpc ($\sim 10^{10}$~ly) \citep{1998riess}. 

Figure~\ref{ladder} illustrates the cosmic distance ladder, beginning with geometric parallax, progressing through Cepheids, and culminating in SNIa measurements at high redshift. Applying this calibrated framework to the Hubble--Lemaître relation yields
\[
H_0 = 73.30 \pm 1.04 \; \mathrm{km\,s^{-1}\,Mpc^{-1}},
\]
which differs by $\sim 5\sigma$ from the value inferred by \emph{Planck} CMB observations under $\Lambda$CDM assumptions \citep{2022riess,2018planck}. This persistent discrepancy, known as the \emph{Hubble tension}, may reflect unaccounted systematic uncertainties or signal new physics beyond the standard cosmological model, and remains a central challenge in contemporary cosmology.


\subsection{Instruments of furture}
The James Webb Space Telescope (JWST), launched in 2021, complements Gaia by extending observations into the infrared, where Cepheids are less affected by dust extinction and metallicity effects. JWST’s sensitivity allows high-precision photometry of Cepheid variables in nearby and distant galaxies, further refining their period--luminosity relation and extending the reach of the cosmic distance ladder deeper into the Hubble flow.

In addition to these missions, major surveys and observatories such as SDSS, Planck, and 2dFGRS have provided critical data for cross-checking and anchoring distance measurements. Future facilities—including the Vera C. Rubin Observatory, DESI, 4MOST, Euclid, SKA, and LISA—promise to deliver even more precise observations, offering independent calibration pathways for standard candles and the local expansion rate.

Beyond photometry and astrometry, the advent of multimessenger astronomy—through the detection of cosmic rays, neutrinos, and gravitational waves—introduces alternative distance indicators, such as standard sirens, which can independently validate Cepheid- and SNIa-based distances. 

Collectively, these technological and observational advances mark the golden era of modern astronomy. By providing highly accurate distances,  they enable precise calibration of the Leavitt Law, strengthening the foundation of the cosmic distance ladder and advancing our ability to determine the Hubble constant with reduced systematic uncertainties.
